% Problem record op_1fc674a1578b1e05. This TeX file is derived from
% database/problems_json/op_1fc674a1578b1e05.json by scripts/sync-tex.mjs; edit the JSON record
% and rerun the script rather than editing this file.
\section{Low-energy Holevo additivity for a squeezed thermal attenuator}
\label{sec:op_1fc674a1578b1e05}

\paragraph{Problem.}
Is energy-constrained Holevo information additive for a squeezed thermal attenuator below its water-filling threshold?
Fix $0<\eta<1$, $b\geq0$, and $r\in\mathbb R\setminus\{0\}$. The channel $\Phi_{\eta,b,r}$ mixes an input mode with an independent centered squeezed thermal mode on a beam splitter of transmissivity $\eta$. Each use has a fresh environment. With $[q,p]=i$, $a=(q+ip)/\sqrt2$, and $R=(q,p)^T$, use $V_{jk}=\langle\{R_j-\langle R_j\rangle,R_k-\langle R_k\rangle\}\rangle/2$. The environment covariance is $(b+1/2)\operatorname{diag}(e^{2r},e^{-2r})$.
Choose an energy $0<E<E_{\mathrm{thr}}$, where
\begin{equation}
 E_{\mathrm{thr}}=\frac{e^{2|r|}-1}{2}
 +\frac{(1-\eta)(b+1/2)}{\eta}\sinh(2|r|).
 \label{eq:ps-threshold}
\end{equation}
For an integer $m\geq1$, set $H_m=\sum_{j=1}^m a_j^\dagger a_j$. Define the unrestricted Holevo information by
\begin{equation}
 \begin{aligned}
 \chi_m(E)&=\sup_{\{p_x,\rho_x\}:\,\operatorname{Tr}\bar\rho H_m\leq mE}
 \left[S(\Phi_{\eta,b,r}^{\otimes m}(\bar\rho))
 -\sum_xp_xS(\Phi_{\eta,b,r}^{\otimes m}(\rho_x))\right],\\
 \bar\rho&=\sum_xp_x\rho_x,\qquad S(\rho)=-\operatorname{Tr}\rho\log_2\rho.
 \end{aligned}
 \label{eq:ps-holevo}
\end{equation}
The states in Eq.~\eqref{eq:ps-holevo} may be entangled across uses. Continuous ensembles are allowed, with sums replaced by integrals. Does $\chi_m(E)=m\chi_1(E)$ hold for every $m$ throughout the regime in Eq.~\eqref{eq:ps-threshold}?

\subsection*{Status}

Unsolved

\subsection*{Source}

This is the squeezed-attenuator specialization of the low-energy additivity gap discussed by Sch\"afer et al., around Corollary~2 and their conclusion \sourcecite{ref:ps-equivalence}{SKG+13}. Ji's version~2 explicitly retains the multiple-use gap in its Scope paragraph and Supplemental Material, Sec.~13 \sourcecite{ref:ps-ji}{Ji26}.

\subsection*{Progress}

\begin{itemize}
  \item For $E\geq E_{\mathrm{thr}}$, a Gaussian distribution of displaced squeezed-vacuum letters attains the thermal maximum-output-entropy bound minus the minimum output entropy. The resulting value per use is $g_2(\eta E+(1-\eta)[(b+1/2)\cosh(2r)-1/2])-g_2((1-\eta)b)$, where $g_2(x)=(x+1)\log_2(x+1)-x\log_2x$ and $0\log_2 0=0$. Corollary~2 of Sch\"afer et al. gives the threshold construction; the subsequent Gaussian optimizer theorem supplies the unrestricted minimum-entropy result \sourcecite{ref:ps-equivalence}{SKG+13}, \sourcecite{ref:ps-optimizer}{GHG15}.

  \item Ji's September 2026 preprint, version~2, Theorem~2, claims $\chi_1(E)=\chi_1^{\mathrm G}(E)$ at every energy. Here the Gaussian restriction uses displaced pure Gaussian letters with Gaussian modulation. The theorem is explicitly one-use and does not establish the equality for $m>1$ \sourcecite{ref:ps-ji}{Ji26}.

  \item The entanglement-breaking regime $b\geq\eta/(1-\eta)$ is additive at every energy. Squeezing equivalence preserves entanglement breaking, and the energy-constrained additivity theorem applies to the resulting Gaussian channel. See Sec.~3, Eqs.~(20)--(24), of Giovannetti, Holevo, and Garc\'ia-Patr\'on \sourcecite{ref:ps-optimizer}{GHG15}.
\end{itemize}

\subsection*{References}

\begin{enumerate}[label={},leftmargin=4.8em,labelsep=0.5em]
  \footnotesize
  \item[\textup{[SKG+13]}]\label{ref:ps-equivalence}
  J. Sch\"afer, E. Karpov, R. Garc\'ia-Patr\'on, O. V. Pilyavets, and N. J. Cerf, ``Equivalence Relations for the Classical Capacity of Single-Mode Gaussian Quantum Channels,'' \emph{Physical Review Letters} \textbf{111}, 030503 (2013). \href{https://doi.org/10.1103/PhysRevLett.111.030503}{doi:10.1103/PhysRevLett.111.030503}; \href{https://arxiv.org/abs/1303.4939}{arXiv:1303.4939}.

  \item[\textup{[GHG15]}]\label{ref:ps-optimizer}
  V. Giovannetti, A. S. Holevo, and R. Garc\'ia-Patr\'on, ``A Solution of Gaussian Optimizer Conjecture for Quantum Channels,'' \emph{Communications in Mathematical Physics} \textbf{334}, 1553--1571 (2015). \href{https://doi.org/10.1007/s00220-014-2150-6}{doi:10.1007/s00220-014-2150-6}; \href{https://arxiv.org/abs/1312.2251}{arXiv:1312.2251}.

  \item[\textup{[Ji26]}]\label{ref:ps-ji}
  S.-W. Ji, ``Gaussian Optimality of Energy-Constrained One-Shot Communication through Single-Mode Bosonic Gaussian Channels,'' arXiv preprint, version~2 (8 September 2026). \href{https://arxiv.org/abs/2608.17239v2}{arXiv:2608.17239v2}.
\end{enumerate}

\subsection*{Comment}

The question allows arbitrary block ensembles and fixes the average photon budget per use. The unsettled regime is below threshold and outside the entanglement-breaking region. One-use Gaussian optimality does not settle regularization. The cited 2026 claim is a preprint result. Squeezing the channel into phase-insensitive form changes the input cost, so phase-insensitive capacity additivity alone does not resolve this question.

\subsection*{Field}

Quantum Communication

\subsection*{Topic}

Classical capacity; Bosonic channels; Additivity and regularization

\subsection*{ID}

\texttt{op\_1fc674a1578b1e05}
